\documentclass[conference]{IEEEtran}

\usepackage{amsmath}
\usepackage{amssymb}
\usepackage{booktabs}
\usepackage{url}
\usepackage[hidelinks]{hyperref} % keep last

\newcommand{\prismloc}{\textsc{prism\_loc}}

\begin{document}

\title{PRISM-Loc: Three LiDAR Localization Backends Behind One
ROS~2 Contract, with Middleware-Free Estimator Cores}

\author{\IEEEauthorblockN{Jung Mo Kang}
\IEEEauthorblockA{kangjmo91@gmail.com}}

\maketitle

\begin{abstract}
Map-based LiDAR localization for mobile robots is solved in fragments:
a 2D scanner calls for an AMCL-style particle filter over an occupancy
grid, a 3D LiDAR for an NDT-based tracker over a point-cloud map, and
an inertial- or GNSS-aided platform for a Kalman fusion node---each
with its own topics, parameters, initialization ritual, and recovery
behavior. We present \prismloc{}, an open-source ROS~2 localization
stack that places three known estimator families behind one node
contract: \texttt{laser2d}, a likelihood-field Monte Carlo localizer;
\texttt{ndt3d}, an NDT-MCL backend that reuses the same particle filter
with a Gaussian-voxel measurement model; and \texttt{fusion3d}, a
15-state error-state Kalman filter fusing LiDAR, IMU, and RTK-GNSS.
Every backend publishes the same \texttt{map}$\to$\texttt{odom}
transform and pose with covariance and accepts the same
\texttt{/initialpose} seeding, so the sensor suite selects a backend at
launch time rather than a separate software stack. All filter
mathematics reside in middleware-free C++17 cores that depend only on
Eigen and are exercised by deterministic, seeded unit and integration
tests with no ROS runtime; the \texttt{laser2d} backend additionally
recovers from the kidnapped-robot condition by deterministic
branch-and-bound scan matching. We claim no new estimator theory: the
contribution is architectural, delivered as a usability-focused
open-source release.
\end{abstract}

\section{Introduction}\label{sec:intro}

Localizing a mobile robot against a prior map is, on paper, a solved
problem, but in practice it is solved in fragments. A team fielding a 2D
LiDAR reaches for \texttt{nav2\_amcl}, a particle filter that tracks a
planar pose against an occupancy grid and nothing else. A team with a 3D
LiDAR reaches instead for an \texttt{hdl\_localization}-style tracker
built around 3D NDT registration. A team that must fuse inertial and
satellite measurements reaches for a \texttt{robot\_localization}-style
node~\cite{moore2014generalized}, which composes an EKF over odometry,
IMU, and GNSS but treats the map-relative LiDAR measurement as someone
else's problem. Each of these tools carries its own topic interface,
parameter vocabulary, initialization ritual, and recovery story, so a
practitioner who moves from a 2D platform to a 3D one, or who adds a
GNSS receiver, re-integrates from scratch. The ROS~2 maintainers'
localization survey~\cite{macenski2023survey} names exactly this gap
and sketches, as unbuilt future work, a modular framework of swappable
2D and 3D localization components behind a common interface
(Section~\ref{sec:related} details this and the surrounding prior
work). The cost of the gap is not merely inconvenience: the Nav2
marathon deployment~\cite{macenski2020marathon} attributes a leading
share of its autonomy-recovery triggers to loss of localization
confidence, so how a stack initializes, degrades, and recovers is a
first-class deployment concern rather than a detail.

\prismloc{} is an engineering response to that gap. It packages three
map-based LiDAR localization backends behind one ROS~2 node contract:
\texttt{laser2d}, a 2D likelihood-field Monte Carlo localizer against an
occupancy grid; \texttt{ndt3d}, a 3D NDT-MCL backend that reuses the
same particle filter with a voxel-Gaussian measurement model over a
prior point cloud, sampling the planar $(x, y, \theta)$ pose; and
\texttt{fusion3d}, a 15-state error-state Kalman filter (ESKF) that fuses
3D LiDAR, IMU, and RTK-GNSS with IMU-bias estimation. All three present the
identical interface --- a \texttt{map}$\to$\texttt{odom} transform, a pose
with $6\times6$ covariance, and \texttt{/initialpose} seeding --- so
that choosing between a 2D scan, a 3D cloud, and an inertial/GNSS-aided
configuration is a launch-time decision rather than a change of stack.
The two MCL backends share a single pure-C++17 estimator core (motion
model, particle filter, KLD-adaptive resampling, pose estimate), and the
fusion backend is a second such core; in both cases the algorithmic code
depends only on Eigen and is unit-tested with no ROS or PCL present.
When \texttt{laser2d} is started with no initial pose, it recovers one
from a single scan by deterministic branch-and-bound (BBS) correlative
scan matching over the occupancy grid, and re-seeds the particle filter
from that result. We claim no new estimator theory: likelihood-field
MCL, NDT-MCL, ESKF fusion, and BBS scan matching are all established.
The contribution is architectural --- three known estimator families
placed behind one contract, factored so their cores are deterministically
testable off-robot --- and it is delivered as usable open-source
software.

Concretely, this paper contributes:
\begin{enumerate}
  \item A \textbf{unified, pluggable localization architecture} in which
    three estimator families (2D likelihood-field MCL, 3D NDT-MCL, and
    LiDAR+IMU+RTK-GNSS ESKF fusion) present one ROS~2
    \texttt{map}$\to$\texttt{odom} contract, so the sensor suite selects
    a backend rather than a separate stack.
  \item \textbf{ROS-free estimator cores} that hold all filter
    mathematics behind an Eigen-only, PCL-free, rclcpp-free boundary,
    making every algorithmic claim exercisable by deterministic,
    seeded-RNG unit and integration tests that run without a middleware
    runtime.
  \item A \textbf{deterministic global relocalization} path for
    \texttt{laser2d}: branch-and-bound correlative scan matching solves
    the kidnapped-robot problem from one scan and re-seeds the MCL
    posterior, replacing stochastic random-injection recovery with a
    repeatable search.
  \item An \textbf{open-source release engineered for usability},
    including fail-fast configuration validation, startup watchdog
    diagnostics that name the input that has gone silent, and a
    single-source parameter reference kept next to the code that reads
    it.
\end{enumerate}

Section~\ref{sec:related} situates \prismloc{} against prior
localization and fusion work. Section~\ref{sec:architecture} describes
the shared contract and the core/ROS boundary; Section~\ref{sec:cores}
details the estimator cores and Section~\ref{sec:ros} the ROS~2 nodes.
Section~\ref{sec:validation} presents the ROS-free testing and CI
strategy that stands in for on-robot benchmarks in this
software-description paper, and Sections~\ref{sec:limitations}
and~\ref{sec:conclusion} discuss limitations and conclude.

\section{Related Work}\label{sec:related}

\textbf{Monte Carlo localization.}
Particle-filter localization was introduced by Dellaert et
al.~\cite{dellaert1999montecarlo} and Fox et
al.~\cite{fox1999efficient}, who represent the pose posterior as a
weighted sample set updated by motion-model prediction,
sensor-likelihood weighting, and resampling, and who demonstrated
global localization from a uniform prior on real robots. Thrun et
al.~\cite{thrun2001robust} later hardened MCL against accurate sensors
and the kidnapped-robot problem via Mixture-MCL's dual sampler. The
particle-filter core in \prismloc{} is a direct descendant of
this line of work and of its \texttt{nav2\_amcl} lineage; we contribute
no new filter theory. What differs is factoring: the same
predict--weight--resample core is shared, unchanged, between a 2D
likelihood-field backend and a 3D NDT backend, and recovery from
localization failure is delegated to a deterministic branch-and-bound
search (Section~\ref{sec:global}) rather than the stochastic
random-injection or dual-sampling mechanisms these papers employ.

\textbf{NDT representations and NDT-based measurement models.}
Biber and Stra{\ss}er~\cite{biber2003ndt} introduced the Normal
Distributions Transform, modeling a scan as per-cell Gaussians and
reducing registration to Newton optimization of a smooth score;
Magnusson et al.~\cite{magnusson2007scan} generalized the
representation to 3D voxel grids for field deployment on mining
vehicles. Both use NDT for deterministic scan-to-scan or scan-to-map
registration that requires a reasonable initial guess.
\prismloc's \texttt{ndt3d} backend adopts the standard
Gaussian-voxel representation but repurposes the NDT score as a
per-particle likelihood inside the shared MCL core (i.e., NDT-MCL
rather than NDT registration), which trades the quadratic convergence
of Newton-based alignment for the multi-hypothesis robustness of the
particle filter, and lets the 2D and 3D backends differ only in their
measurement model.

\textbf{LiDAR-inertial and GNSS fusion.}
Sol\`a's tutorial~\cite{sola2017quaternion} is the mathematical basis
of \prismloc's fusion backend: the nominal/error-state
decomposition, local angular-error parameterization, and
predict--correct--inject--reset cycle are implemented essentially as
derived there. On the systems side, Moore and
Stouch's \texttt{robot\_localization}~\cite{moore2014generalized}
established the pattern of a generalized, configurable EKF/UKF fusion
node for ROS, and Wan et al.~\cite{wan2018robust} demonstrated an
industrial ESKF fusing GNSS-RTK, LiDAR, and IMU at centimeter accuracy
for the Baidu Apollo fleet, albeit as a closed system tied to
proprietary intensity/altitude maps. Tightly-coupled LiDAR-inertial
odometry such as FAST-LIO~\cite{xu2021fastlio}, and applied 3D
map-based tracking systems such as Koide et
al.~\cite{koide2019portable}, are complementary: they produce
ego-motion and maps, whereas \prismloc{} localizes within a
prior map. Relative to these, our \texttt{fusion3d} backend claims no
estimator novelty; it packages a textbook ESKF as a ROS-free core
behind the same ROS~2 contract as the two MCL backends, so that
odometry-only, LiDAR-only, and GNSS-aided configurations are
deployment choices rather than separate stacks.

\textbf{Global (kidnapped-robot) localization.}
Olson~\cite{olson2009correlative} formulated 2D scan matching as
exhaustive correlative search over a rasterized log-likelihood table
with coarse-to-fine multi-resolution pruning, and Hess et
al.~\cite{hess2016realtime} bounded that search with branch-and-bound
over a precomputed grid pyramid to run loop closure in real time
inside Cartographer. \prismloc{} borrows exactly this
branch-and-bound-over-likelihood-pyramid primitive but applies it in a
narrower setting: one-shot relocalization against a known static map,
with no submaps or pose graph to maintain. This replaces the heuristic
recovery mechanisms of classical
MCL~\cite{fox1999efficient,thrun2001robust} with a deterministic
search whose result re-seeds the particle filter.

\textbf{Robot software systems.}
Nav2~\cite{macenski2020marathon} treats localization as an external,
swappable component (stock AMCL or SLAM Toolbox), and its marathon
deployment attributes a leading share of recovery triggers to
localization-confidence loss in feature-poor corridors. The ROS~2
maintainers' survey~\cite{macenski2023survey} makes the gap explicit:
it catalogs AMCL's shortcomings and sketches, as unbuilt future work,
a modular framework for composing 2D and 3D localization systems from
swappable components. As framed in Section~\ref{sec:intro}, \prismloc{}
answers with engineering rather than a new algorithm: three known
estimator families (likelihood-field MCL, NDT-MCL, ESKF fusion) are
placed behind one ROS~2 topic and TF contract, with each estimator
implemented as a ROS-free library that is unit-testable without a
middleware runtime --- a property that \texttt{robot\_localization}'s
otherwise modular design~\cite{moore2014generalized} did not
prioritize. The accuracy evaluation planned in
Section~\ref{sec:limitations} will follow the trajectory-error
conventions (ATE/RPE) standardized by the TUM RGB-D
benchmark~\cite{sturm2012benchmark}, applied to per-backend
\texttt{map}$\to$\texttt{odom} estimates rather than visual SLAM
output.

\section{System Architecture}\label{sec:architecture}

\prismloc{} is organized as two layers realized by four packages
(Fig.~\ref{fig:layers}). The lower layer holds the estimator cores:
\texttt{prism\_loc\_core}, which implements the Monte Carlo
localization (MCL) family, and \texttt{prism\_loc\_fusion}, which
implements the error-state Kalman filter (ESKF). Both are plain
C++17 libraries whose only third-party dependency is Eigen; each
\texttt{CMakeLists.txt} calls \texttt{find\_package(Eigen3 REQUIRED)}
and links \texttt{Eigen3::Eigen} alone, with no reference to
\texttt{rclcpp}, \texttt{tf2}, or PCL. The upper layer holds the ROS~2
node packages: \texttt{prism\_loc} (the \texttt{laser2d} and
\texttt{ndt3d} backends) and \texttt{prism\_loc\_fusion\_ros} (the
\texttt{fusion3d} backend). These packages, and only these, depend on
\texttt{rclcpp}, the \texttt{tf2} stack, the ROS message packages, and
PCL. Each core is consumed by its node package through the ordinary
ament \texttt{find\_package} mechanism, so the dependency direction is
strictly upward: the cores never link against the middleware.

\begin{figure}[t]
\centering
\footnotesize
\setlength{\tabcolsep}{4pt}
\begin{tabular}{|l|}
\hline
\textbf{Estimator cores} --- C++17 + Eigen, no ROS/PCL, gtest \\
\hline
\texttt{prism\_loc\_core} (MCL):\\
\quad OdometryMotionModel $\rightarrow$ ParticleFilter\\
\quad\quad $\vdash$ Laser2DLikelihoodField\\
\quad\quad $\vdash$ Ndt3DModel\\
\quad $\rightarrow$ Pose2D $+$ $6{\times}6$ covariance\\
\texttt{prism\_loc\_fusion} (ESKF):\\
\quad IMU predict $\rightarrow$ 15-state ESKF\\
\quad\quad $\dashv$ NDT 6-DoF pose update\\
\quad\quad $\dashv$ RTK-GNSS position update (WGS84 $\rightarrow$ ENU)\\
\hline
\multicolumn{1}{c}{}\\[-6pt]
\hline
\textbf{ROS~2 nodes} --- rclcpp / tf2 / PCL only here \\
\hline
\texttt{prism\_loc}\hfill \texttt{backend = laser2d | ndt3d}\\
\texttt{prism\_loc\_fusion\_ros}\hfill \texttt{backend = fusion3d}\\
\quad adapters: laser / pointcloud / ndt\_registration\\
\hline
in:~\texttt{/scan|/points|/imu|/gnss}, map, \texttt{/initialpose}, TF\\
out:~\texttt{/tf} (\texttt{map}$\rightarrow$\texttt{odom}), \texttt{\char`~/pose},
\texttt{\char`~/particle\_cloud|\char`~/odometry}\\
\hline
\end{tabular}
\caption{The two-layer split. Estimator cores (bottom) are
middleware-free libraries; the ROS~2 nodes (top) select a backend at
runtime and bridge messages to core types through adapters. The
dependency arrow points only upward.}
\label{fig:layers}
\end{figure}

\textbf{Why the split.} Isolating the estimators from the middleware
buys three properties. First, unit-testability without a running
graph: each core builds and tests with a plain system toolchain and
gtest, so its numerical behavior is exercised with no ROS installation,
no node, and no message plumbing. The \texttt{CMakeLists.txt} of each
core wires tests in dual mode --- through ament when a ROS workspace is
present, and through bare CTest plus system gtest otherwise --- and this
ROS-free path is the one continuous integration builds. The detailed
testing story is deferred to Section~\ref{sec:validation}. Second,
determinism: all sampling routes through a single seeded
\texttt{std::mt19937\_64} wrapper (default seed \texttt{42}), so a
motion-prediction, weighting, or resampling step is reproducible bit
for bit across runs, which is what makes the core tests assertable
rather than merely statistical. Third, one output contract shared by
all three backends. Regardless of backend, a node publishes the
\texttt{map}$\rightarrow$\texttt{odom} transform on \texttt{/tf}, a
\texttt{PoseWithCovarianceStamped} on \texttt{\char`~/pose}, and
accepts a \texttt{PoseWithCovarianceStamped} on \texttt{/initialpose}
for seeding; the MCL nodes additionally publish a particle cloud while
the fusion node publishes an \texttt{Odometry} message carrying pose
and velocity. Because the contract is fixed, choosing between the
likelihood-field MCL, NDT-MCL, and ESKF estimators is a runtime
decision: \texttt{prism\_loc} reads a \texttt{backend} string parameter
(\texttt{"laser2d"} by default) and branches on it when constructing
its measurement model and when applying the correction step, and
\texttt{prism\_loc\_fusion\_ros} supplies the \texttt{fusion3d}
estimator behind the identical topic and TF interface. As stated in
Section~\ref{sec:related}, this factoring --- three known estimator
families behind one ROS~2 contract --- is the architectural
contribution; the estimators themselves are textbook.

\textbf{Core interfaces.} The MCL core expresses the shared
predict--weight--resample loop against a small set of value types.
A \texttt{Pose2D} carries the planar pose $(x, y, \theta)$; a \texttt{Particle}
pairs a pose with a weight; and the two backends differ only through a
\texttt{MeasurementModel} abstraction whose single virtual method
returns the log-likelihood of an observation given a candidate pose.
\texttt{Laser2DLikelihoodField} (a 2D scan against an occupancy-grid
likelihood field) and \texttt{Ndt3DModel} (a 3D cloud against a
Gaussian-voxel NDT map) are its two implementations, and the same
\texttt{ParticleFilter} and \texttt{OdometryMotionModel} drive both.
The fusion core exposes an \texttt{Eskf} class over a 15-dimensional
error state with three entry points --- an IMU-driven
\texttt{predict}, a 6-DoF \texttt{updatePose} fed by NDT registration,
and a 3-DoF \texttt{updatePosition} fed by an absolute position fix ---
alongside a geodetic helper that converts WGS84 latitude/longitude/
altitude to a local ENU frame. Section~\ref{sec:cores} details these
estimators.

\textbf{Adapters.} The node layer bridges ROS message types to the
middleware-free core types through thin adapters, keeping the cores
ignorant of message definitions. In \texttt{prism\_loc}, a laser
adapter converts a \texttt{LaserScan} to the core's \texttt{LaserScan2D}
(honoring range overrides) and an \texttt{OccupancyGrid} to the core
grid map, and a pointcloud adapter converts a \texttt{PointCloud2}, and
a \texttt{.pcd} map file, into the \texttt{std::vector<Eigen::Vector3d>}
the NDT map ingests; both also convert quaternion-bearing
\texttt{Transform} and \texttt{Pose} messages to \texttt{Pose2D} by
extracting yaw. In \texttt{prism\_loc\_fusion\_ros}, an
\texttt{ndt\_registration} adapter wraps PCL's NDT to align an incoming
cloud against the prior-map cloud and returns the resulting
\texttt{Isometry3d} pose, which becomes the ESKF's 6-DoF measurement.
PCL therefore appears only inside these adapters, never in the cores.
The node-layer responsibilities --- TF lookups, topic wiring, and
runtime diagnostics --- are described in Section~\ref{sec:ros}.

\section{Estimator Cores}\label{sec:cores}

This section describes the three estimator families that \prismloc\ places
behind a single ROS~2 contract. All three are implemented as ROS-free C++
libraries: the two Monte Carlo localization (MCL) backends and the global
relocalizer live in \texttt{prism\_loc\_core}, and the fusion filter lives in
\texttt{prism\_loc\_fusion}. Consistent with the architectural framing of
Section~\ref{sec:architecture}, the contribution here is factoring, not new
estimation theory: we reuse textbook constructions and report only what the
code implements.

\subsection{Shared Monte Carlo Localization Core}

The two 2D and 3D pose-tracking backends share one particle filter. Following
Thrun-style MCL, the pose posterior at time $t$ is a weighted sample set
$\{(x_t^{[i]}, w_t^{[i]})\}_{i=1}^{N}$ over the planar pose
$x_t = (x, y, \theta)$ (\texttt{Pose2D}), updated by the standard
predict--weight--resample recursion. Only the measurement model differs between
backends, so the entire filter is written against an abstract interface rather
than any specific sensor.

\paragraph{Motion model} \texttt{OdometryMotionModel} implements the
sampling-based odometry model. Each control $u_t$ is decomposed into a first
rotation, a translation, and a second rotation
$(\delta_{r1}, \delta_{t}, \delta_{r2})$ from consecutive odometry poses, and
each particle is advanced by the noise-corrupted increment
\begin{align}
\hat{\delta}_{r1} &= \delta_{r1} - \varepsilon(\alpha_1 \delta_{r1}^2 + \alpha_2 \delta_{t}^2), \nonumber\\
\hat{\delta}_{t}  &= \delta_{t}  - \varepsilon(\alpha_3 \delta_{t}^2 + \alpha_4(\delta_{r1}^2 + \delta_{r2}^2)), \nonumber\\
\hat{\delta}_{r2} &= \delta_{r2} - \varepsilon(\alpha_1 \delta_{r2}^2 + \alpha_2 \delta_{t}^2),
\end{align}
where $\varepsilon(\sigma^2)$ draws zero-mean Gaussian noise with the given
variance and $\alpha_{1\text{--}4}$ (\texttt{MotionParams}) are the usual
rotation/translation noise coefficients. A pure-rotation guard sets
$\delta_{r1}=0$ when the translation is below $10^{-3}$~m to avoid an ill-defined
heading.

\paragraph{Particle filter} \texttt{ParticleFilter} weights each particle by the
measurement model's log-likelihood, subtracting the per-step maximum before
exponentiation for numerical stability, then normalizing. Resampling is
triggered only when the effective sample size
$N_\mathrm{eff} = (\sum_i w^{[i]})^2 / \sum_i (w^{[i]})^2$ drops below
\texttt{resample\_threshold} times $N$, and is KLD-adaptive~\cite{fox2003kld}:
samples are drawn proportional to weight until the count reaches the
Kullback--Leibler bound
\begin{equation}
M_x = \frac{k-1}{2\,\varepsilon_\mathrm{kld}}
      \left(1 - \frac{2}{9(k-1)} + \sqrt{\tfrac{2}{9(k-1)}}\,z_{1-\delta}\right)^{3},
\end{equation}
where $k$ is the number of occupied histogram bins seen so far. Bins are formed
by discretizing $(x,y)$ at \texttt{kld\_bin\_xy} and $\theta$ at
\texttt{kld\_bin\_yaw}; the sample count is clamped to
$[\texttt{min\_particles}, \texttt{max\_particles}]$, and $z_{1-\delta}$
(\texttt{kld\_z}) and $\varepsilon_\mathrm{kld}$ (\texttt{kld\_err}) are the
standard-normal quantile and error bound. The pose estimate is the weighted mean
of $(x,y)$ and the circular mean of $\theta$; the reported $6\times6$ covariance
carries the planar terms and marks $z$, roll, and pitch as unobserved.

\paragraph{Measurement models} A backend supplies a \texttt{MeasurementModel},
whose sole method returns the unnormalized log-likelihood of the current
observation given a candidate pose. Two are provided.
\texttt{Laser2DLikelihoodField} is a Thrun-style likelihood-field model for 2D
scans~\cite{thrun2001robust}: an occupancy grid is preprocessed into a
\texttt{LikelihoodField} that stores, per cell, the Euclidean distance $d$ to the
nearest occupied cell (computed once by a separable distance transform), and each
subsampled beam endpoint contributes
$\log\!\big(z_\mathrm{hit}\,e^{-d^2/2\sigma_\mathrm{hit}^2} + z_\mathrm{rand}/z_\mathrm{max}\big)$,
with \texttt{z\_hit}, \texttt{z\_rand}, \texttt{sigma\_hit}, and
\texttt{max\_beams} as parameters. \texttt{Ndt3DModel} reuses the identical
filter but scores a 3D point cloud against a Normal Distributions Transform
map~\cite{biber2003ndt,magnusson2007scan}. The prior map is voxelized into
per-voxel Gaussians $(\mu_v, \Sigma_v)$ (\texttt{NdtMap}, with small eigenvalues
inflated to a fraction of the largest for numerical conditioning), and the
per-particle likelihood accumulates
$\log\!\big(s(p) + s_\mathrm{floor}\big)$ over cloud points $p$, where
$s(p) = \exp\!\big(-\tfrac{1}{2}(p-\mu_v)^\top \Sigma_v^{-1}(p-\mu_v)\big)$ and
empty voxels score zero. The planar particle state $(x,y,\theta)$ is lifted to a
3D transform at a fixed base height, so the 3D backend estimates the same
$(x,y,\theta)$ as the 2D one. This is the NDT-as-measurement-model (NDT-MCL)
idea: the NDT score usually optimized by Newton registration is instead read out
as a particle likelihood, so the two backends differ only in this class.

\subsection{Branch-and-Bound Global Localization}\label{sec:global}

Recovery from a lost or kidnapped estimate is handled by a deterministic
correlative scan matcher, \texttt{BranchAndBoundMatcher}, following Olson's
multi-resolution correlative search~\cite{olson2009correlative} and the
branch-and-bound formulation of Hess et al.~\cite{hess2016realtime}. At
construction the likelihood field is converted into a probability grid
$p = e^{-d^2/2\sigma_\mathrm{hit}^2}$, and a max-pooling pyramid of
\texttt{max\_depth}$+1$ levels is precomputed, where level $h$ stores the maximum
over a $2^h\times2^h$ cell block. A search then proceeds over discretized poses:
translations within \texttt{linear\_window} and rotations spanning
\texttt{angular\_window} at \texttt{angular\_step} resolution. Coarse nodes are
expanded best-first and pruned whenever their pyramid-derived upper bound cannot
exceed the best exact (leaf) score found so far, which keeps the bound
admissible and the returned pose globally optimal over the discretization. The
result is accepted only when its score reaches \texttt{min\_score\_fraction} of
the number of beams used; an accepted pose re-seeds the particle filter, whereas
a rejected one leaves the current estimate untouched. Unlike the stochastic
random-injection recovery of classical MCL, this search is deterministic and
therefore directly reproducible in a unit test.

\subsection{Error-State Kalman Filter Fusion}

The fusion backend, \texttt{Eskf}, is a 15-state error-state Kalman filter that
follows Sol\`a's quaternion formulation~\cite{sola2017quaternion}. The nominal
state (\texttt{NominalState}) carries position $p$, velocity $v$, orientation
quaternion $q$, accelerometer bias $b_a$, and gyroscope bias $b_g$; the error
state $\delta x \in \mathbb{R}^{15}$ collects
$(\delta p, \delta v, \delta\theta, \delta b_a, \delta b_g)$ with a local
angular error $\delta\theta$. On each IMU sample the nominal state is integrated
in closed form,
\begin{equation}
p \mathrel{+}= v\,\Delta t + \tfrac{1}{2} a_w \Delta t^2,\quad
v \mathrel{+}= a_w \Delta t,\quad
q \leftarrow q \otimes \mathrm{Exp}(\omega\,\Delta t),
\end{equation}
where $a_w = R(q)(a_m - b_a) + g$ and $\omega = \omega_m - b_g$ are the
bias-corrected, gravity-compensated measurements, and $\mathrm{Exp}(\cdot)$ is
the $SO(3)$ exponential (\texttt{so3.hpp}). The error covariance is propagated
with the discrete transition $F$ and additive process noise $Q$ built from the
IMU noise parameters,
\begin{equation}
P \leftarrow F\,P\,F^\top + Q .
\end{equation}
Two measurement updates share the standard Kalman gain
$K = P H^\top (H P H^\top + R)^{-1}$, $\delta x = K\,y$,
$P \leftarrow (I - KH)P$: \texttt{updatePose} applies a 6-DoF NDT pose
observation whose residual stacks the position error and the local orientation
error $\log(q^{-1} q_\mathrm{meas})$, while \texttt{updatePosition} applies a
3-DoF GNSS position observation. GNSS fixes are mapped from geodetic
coordinates into a local ENU frame through a WGS-84 LLA$\rightarrow$ECEF$\rightarrow$ENU
conversion about a fixed datum (\texttt{GeodeticConverter}). After each update
the estimated error is injected into the nominal state --- additively for the
vector components and multiplicatively, $q \leftarrow q \otimes
\mathrm{Exp}(\delta\theta)$, for the orientation --- and the error state is reset
to zero. Because the filter takes plain vectors and quaternions rather than ROS
messages, the entire predict--correct--inject--reset cycle is exercised in ROS-free unit
tests.

\section{ROS 2 Integration and Interface Contract}\label{sec:ros}

The estimator cores of Section~\ref{sec:cores} carry no middleware
dependency; they are made deployable by thin ROS~2 node packages
(\texttt{prism\_loc} for the \texttt{laser2d} and \texttt{ndt3d} MCL
backends, \texttt{prism\_loc\_fusion\_ros} for the \texttt{fusion3d}
ESKF backend) whose job is to translate messages, TF, and parameters
into core calls and back. We treat this boundary as a first-class
interface rather than glue code, following the conventions that
\texttt{robot\_localization}~\cite{moore2014generalized} and
\texttt{nav2\_amcl} within Nav2~\cite{macenski2020marathon}
established for drop-in localization: a REP-105 frame contract, a
\texttt{map}$\to$\texttt{odom} TF output, and \texttt{/initialpose}
seeding. A robot already configured for AMCL can therefore switch
backends by changing a launch file, not its TF tree or its consumers.

\subsection{Frames and \texttt{map}$\to$\texttt{odom} composition}
All three backends honor the REP-105 frame chain
\texttt{map}$\to$\texttt{odom}$\to$\texttt{base\_link}, exposed as the
\texttt{global\_frame}, \texttt{odom\_frame}, and \texttt{base\_frame}
parameters. Each backend estimates the robot body in the map,
$T_{\mathrm{map,base}}$, but---so as not to fight the odometry source
that owns \texttt{odom}$\to$\texttt{base\_link}---it broadcasts only
the correction transform $T_{\mathrm{map,odom}}$. Writing $\oplus$ for
pose composition and $(\cdot)^{-1}$ for inversion, the node computes
\begin{equation}
  T_{\mathrm{map,odom}} \;=\;
  T_{\mathrm{map,base}} \oplus T_{\mathrm{odom,base}}^{-1},
  \label{eq:mapodom}
\end{equation}
where $T_{\mathrm{odom,base}}$ is read from TF at the measurement
stamp. Equation~\eqref{eq:mapodom} guarantees the round-trip identity
$T_{\mathrm{map,odom}} \oplus T_{\mathrm{odom,base}} =
T_{\mathrm{map,base}}$, so downstream consumers recover the intended
body pose; both the planar core
(\texttt{compose}/\texttt{inverse} on $SE(2)$) and the fusion node
(\texttt{map\_base}$\cdot$\texttt{odom\_base}$^{-1}$ on $SE(3)$
Isometries) implement it, and each is checked by a ROS-free unit test
asserting the round trip to $10^{-9}$. If \texttt{odom}$\to$\texttt{base\_link}
is momentarily unavailable, the fusion node degrades gracefully to
publishing \texttt{map}$\to$\texttt{base\_link} directly rather than
emitting a wrong transform. The broadcast stamp is future-dated by
\texttt{transform\_tolerance} (default $0.1$\,s) to bound
extrapolation on the consumer side.

\subsection{Topic contract and QoS}
Table~\ref{tab:contract} lists the per-backend contract. Inputs and
outputs are uniform across backends wherever the sensing allows: every
backend outputs the same \texttt{map}$\to$\texttt{odom} TF and a
\texttt{PoseWithCovarianceStamped} on \texttt{\textasciitilde/pose},
and every backend accepts \texttt{/initialpose}. QoS profiles match
each publisher's expected semantics rather than using defaults: sensor
streams (\texttt{LaserScan}, \texttt{PointCloud2}, \texttt{Imu},
\texttt{NavSatFix}) are subscribed with \texttt{SensorDataQoS}
(best-effort, volatile), while the occupancy-grid map is subscribed
\texttt{transient\_local} and \texttt{reliable} so that the single
latched grid published once by \texttt{map\_server} is delivered even
though the subscriber joins later. A mismatch here is a common and
silent failure mode for new ROS users; fixing the profiles at the
interface removes it.

\begin{table}[t]
\centering
\caption{Per-backend topic contract. All backends additionally
subscribe \texttt{/initialpose} and the TF
\texttt{odom}$\to$\texttt{base\_link}, and broadcast
\texttt{map}$\to$\texttt{odom}.}
\label{tab:contract}
\small
\begin{tabular}{@{}lll@{}}
\toprule
Backend & Inputs & Outputs \\
\midrule
\texttt{laser2d} & \texttt{/scan}, \texttt{/map} & \texttt{\textasciitilde/pose}, \\
 & (\texttt{OccupancyGrid}) & \texttt{\textasciitilde/particle\_cloud} \\[2pt]
\texttt{ndt3d} & \texttt{/points}, \texttt{.pcd} & \texttt{\textasciitilde/pose}, \\
 & (file) & \texttt{\textasciitilde/particle\_cloud} \\[2pt]
\texttt{fusion3d} & \texttt{/points}, \texttt{/imu}, & \texttt{\textasciitilde/pose}, \\
 & \texttt{/gnss} (\texttt{NavSatFix}), \texttt{.pcd} & \texttt{\textasciitilde/odometry} \\
\bottomrule
\end{tabular}
\end{table}

\subsection{Initialization paths}
The interface offers one manual and one automatic seeding path per
backend. Manually, all backends accept a
\texttt{PoseWithCovarianceStamped} on \texttt{/initialpose} (the RViz
``2D Pose Estimate'' tool): the MCL backends spread a Gaussian particle
cloud at the supplied mean and covariance, and the fusion backend loads
the pose as its initial nominal state. Automatically, the two backends
differ because their sensing differs. The \texttt{laser2d} backend can
self-recover with no prior pose by running the branch-and-bound search
of Section~\ref{sec:global} (\texttt{try\_global\_localization}, or the
on-demand \texttt{\textasciitilde/global\_localization} service), which
re-seeds the filter from a single scan. The \texttt{fusion3d} backend
auto-initializes once it has both an attitude estimate---roll and pitch
recovered from the IMU specific-force vector, heading from
\texttt{initial\_yaw}---and a position fix from the first accepted GNSS
message, converted from WGS-84 to a local ENU frame; \texttt{/initialpose}
overrides this when GNSS is unavailable.

\subsection{Parameter surface}
Every tunable is a declared ROS~2 parameter with an in-code default,
documented in a single source of truth (\texttt{PARAMS.md}) generated
from the \texttt{declare\_parameter} calls. The surface deliberately
exposes the core's control knobs rather than hiding them: the
KLD-sampling bounds (\texttt{kld\_err}, \texttt{kld\_z}, and the
histogram bin sizes), the likelihood-field and laser-range parameters
(\texttt{z\_hit}, \texttt{sigma\_hit}, \texttt{laser\_min\_range},
\texttt{laser\_max\_range}), the branch-and-bound search window and
depth (\texttt{bbs\_linear\_window}, \texttt{bbs\_angular\_step},
\texttt{bbs\_max\_depth}, \texttt{bbs\_min\_score\_fraction}), and the
ESKF process noise and GNSS gates (\texttt{sigma\_acc},
\texttt{gnss\_min\_status}, \texttt{gnss\_max\_pos\_cov}). Because the
cores are configured purely through plain structs, each parameter maps
to exactly one field, and the shipped YAML files override only the
subset a deployment typically changes.

\subsection{Runtime diagnostics}
A localization node that silently does nothing is the dominant
first-run frustration; this release invests in making every stalled
precondition observable and self-explaining. Four mechanisms cooperate.
(i)~\emph{Fail-fast on a missing mandatory map}: the \texttt{ndt3d} and
\texttt{fusion3d} nodes abort construction with an actionable
\texttt{std::runtime\_error} if \texttt{map\_pcd\_path} is empty or
fails to load, rather than idling forever. (ii)~\emph{Throttled
waiting-state warnings}: whenever a callback cannot yet act---no map,
no initial pose, no filter initialization---it emits a rate-limited
\texttt{WARN} that names both the unmet precondition and the remedy
(e.g.\ ``waiting for an initial pose---publish to \texttt{/initialpose}
\ldots or enable \texttt{try\_global\_localization}''). (iii)~A
\emph{$10$\,s startup watchdog} reports, for each mandatory
input that has produced no message, the topic name and its live
publisher count (via \texttt{count\_publishers}), disambiguating ``no
publisher'' from ``publisher present but QoS-incompatible''; it cancels
itself once all inputs are seen. (iv)~\emph{Quantified rejection
logging}: the GNSS gate logs each rejected fix with the measured value
against the threshold it violated---status below \texttt{gnss\_min\_status},
NaN coordinates, or position variance above \texttt{gnss\_max\_pos\_cov}
---and the NDT update likewise reports non-convergence and fitness
versus \texttt{ndt\_max\_fitness}. These diagnostics carry no algorithmic
weight; they exist so that a misconfiguration is diagnosable from the
console log alone, consistent with the operational transparency Nav2
found necessary for long-term autonomy~\cite{macenski2020marathon}.

\section{Implementation and Validation}\label{sec:validation}

\prismloc{} is validated at the level of \emph{correctness}, not
field accuracy. Because every estimator is a ROS-free library
(Section~\ref{sec:cores}), each can be exercised by an ordinary
C++ unit test with no middleware runtime, no simulator, and no
recorded bags. The suite comprises 44 GoogleTest cases across four
packages: 25 in the shared \texttt{prism\_loc\_core} (the two MCL
backends and global recovery), 13 in \texttt{prism\_loc\_fusion} (the
ESKF), 4 in the \texttt{prism\_loc} ROS node's adapter layer, and 2 in
\texttt{prism\_loc\_fusion\_ros}.

\textbf{What the cores are pinned to.}
The core tests target the components most likely to regress silently.
The odometry motion model is checked both deterministically
(zero-command and pure-translation cases) and statistically (the mean
of noisy samples approximates the commanded motion). Measurement models
are checked for the monotonicity property they must satisfy to be usable
as a likelihood: the ground-truth pose must out-score a displaced one,
for both the 2D likelihood field and the 3D NDT voxel score. The
particle-filter mechanics---low-variance resampling concentrating on a
heavy particle, and the weighted-mean estimate---are tested directly, as
are the occupancy-grid and NDT-map primitives (world/map round-trips,
out-of-bounds handling, likelihood-field distances, and voxel-mean
recovery). Global recovery is exercised by having branch-and-bound
recover the correct global pose from a single scan---the kidnapped-robot
scenario---and reject an empty scan. Two integration tests close the
loop end to end, driving the full predict--weight--resample cycle until
the estimate converges to ground truth for the \texttt{laser2d} and
\texttt{ndt3d} backends respectively.

\textbf{Fusion and math.}
The ESKF is pinned for both convergence and consistency: a static IMU
leaves the state at rest, constant acceleration reproduces the expected
kinematics, and position and pose updates pull the state toward the
measurement; an integration test tracks a moving target from a
deliberately wrong initial guess. Supporting identities are verified
independently---$SO(3)$ exponential/logarithm round-trips and the
skew/cross-product relation, geodetic ENU round-trips against a datum,
and the \texttt{map}$\to$\texttt{odom} TF composition arithmetic used to
publish the correction transform.

\textbf{Determinism.}
Every stochastic path draws from a single seedable generator exposed as
the ROS parameter \texttt{seed} (default \texttt{42}). The RNG tests
assert that equal seeds reproduce identical sequences and that distinct
seeds differ, so the statistical assertions above are reproducible run
to run rather than flaky.

\textbf{Continuous integration.}
Two CI jobs guard the contract from both sides. A \emph{core} job builds
\texttt{prism\_loc\_core} and \texttt{prism\_loc\_fusion} with plain
CMake and runs \texttt{ctest}, requiring only a compiler and Eigen---no
ROS---which keeps the cores honestly middleware-independent. A separate
\emph{ROS} job builds inside a \texttt{ros:humble} container, resolves
dependencies with \texttt{rosdep}, and runs \texttt{colcon build} and
\texttt{colcon test} across all four packages, so the ROS~2 wrapping is
tested against a real distribution.

\textbf{Scope.}
This validation establishes that the estimators are individually correct
and that all three sit behind one working ROS~2 contract; it is
deliberately not an accuracy benchmark. The planned quantitative field
evaluation is discussed with the other roadmap items in
Section~\ref{sec:limitations}.

\section{Limitations and Future Work}\label{sec:limitations}

\prismloc{} is an architecture paper, and its scope is deliberately
bounded. We state the current boundaries plainly so that they read as a
roadmap rather than as hidden assumptions.

\textbf{Planar state in the MCL backends.} Both particle-filter
backends estimate a planar pose $(x, y, \theta)$: the shared core
represents each particle as an $SE(2)$ element, and the \texttt{ndt3d}
backend evaluates its per-particle likelihood against a full 3D
Gaussian-voxel map from that planar pose at a fixed base height. This
suits wheeled robots on locally flat floors but does not track roll,
pitch, or vertical motion. Full $SE(3)$ state exists only in the
\texttt{fusion3d} backend, whose $15$-dimensional error state carries
position, velocity, orientation, and IMU biases. Extending the MCL core
to a $6$-DoF particle state, or coupling the planar filter to an
external attitude source, is future work.

\textbf{2D-only global relocalization.} The branch-and-bound search
used for kidnapped-robot recovery (Section~\ref{sec:global}) operates on
a rasterized 2D occupancy pyramid and returns an $SE(2)$ pose. There is
no analogous 3D global localization for the \texttt{ndt3d} or
\texttt{fusion3d} backends; those currently rely on an operator-supplied
initial pose (via \texttt{/initialpose}) and local tracking. A 3D
branch-and-bound or learned place-recognition front end is on the
roadmap.

\textbf{No online sensor time-offset estimation.} The
\texttt{fusion3d} ESKF propagates on IMU samples and corrects with
LiDAR-NDT and GNSS measurements at their stamped times, but it does not
estimate the temporal offset between the LiDAR/GNSS clocks and the IMU
clock, nor does it self-calibrate the LiDAR--IMU extrinsic; both are
treated as known inputs. Jointly estimating time offset and extrinsics
would improve robustness under poorly synchronized hardware.

\textbf{Loosely-coupled LiDAR update.} In \texttt{fusion3d}, PCL NDT
registers each live cloud to the prior map to produce a $6$-DoF pose
\emph{measurement}, which the ESKF then fuses as a loosely-coupled
update. This keeps the ROS-free filter core independent of the
registration engine, but it forgoes the accuracy of tightly-coupled
formulations that fold raw scan residuals directly into the filter
Jacobian, as in FAST-LIO~\cite{xu2021fastlio}. The loose coupling is an
architectural choice favoring modularity and testability over
last-percent accuracy.

\textbf{Single static map.} All three backends assume one prior map of
a static world, loaded once at startup; they neither update the map
online nor reason about moving obstacles or long-term scene change.
Map maintenance and change detection are out of scope for the present
system.

\textbf{No field-accuracy evaluation yet.} The validation reported in
Section~\ref{sec:validation} is deterministic, ROS-free unit and
integration testing plus continuous integration; it establishes that
each estimator core behaves as specified, but it is not a quantitative
accuracy study. We have not yet measured trajectory error on physical
or public datasets. Planned work is an evaluation reporting absolute and
relative trajectory error (ATE/RPE) per the TUM RGB-D
conventions~\cite{sturm2012benchmark} on public datasets, applied to the
per-backend \texttt{map}$\to$\texttt{odom} estimates. Until then we make
no accuracy claims.

\section{Conclusion}\label{sec:conclusion}

We presented \prismloc, a LiDAR localization stack whose contribution is
architectural rather than algorithmic: three known estimator families
--- likelihood-field MCL, NDT-MCL, and an error-state Kalman filter that
fuses LiDAR, IMU, and RTK-GNSS --- are placed behind a single ROS~2
topic and TF contract, so that a 2D LiDAR, a 3D LiDAR, or a
sensor-fused configuration is a deployment choice rather than a separate
software stack. Each estimator is implemented as a pure C++ core with no
dependence on the middleware runtime, making the filters unit-testable
without ROS and letting the two MCL backends share one
predict--weight--resample core that differs only in its measurement
model. We claim no new filter theory; the novelty is the factoring that
makes these families interchangeable and independently testable behind a
familiar interface. \prismloc{} is released under the Apache-2.0 license
at \url{https://github.com/kjungmo/prism_loc}.


\bibliographystyle{IEEEtran}
\bibliography{references}

\end{document}
